\section{A common absorbing set for each class}\label{sec:permanence}

Fix a positive stoichiometric class $P$ and its representative $c$, and
retain the scale $h_P$ and the functions $F_h$ from
Section~\ref{sec:trapping}. Lemma~\ref{lem:trapping} gives invariance
at each scale and finite-time entry whenever a trajectory stays in the
corresponding box. Proposition~\ref{prop:plateau} connects these two
properties through
\begin{equation}\label{eq:scale-containment}
 K_h\subset\operatorname{int}B_{\sqrt h}
 \qquad(0<h\leq h_P).
\end{equation}
We now use this relation to replace an initial-point-dependent trapping
set by one fixed set for $P$.

\begin{proof}[Proof of Theorem~\ref{thm:main}]
Let $x^0\in P$. Since $x^0$ is positive, we may choose
$0<h_0\leq h_P$ so small that
\[
 \norm{\frac{\log x^0}{\log h_0}}_\infty<H,
\]
where the logarithm is componentwise and $H>0$ is the fixed separation
constant from Section~\ref{sec:trapping}. In particular, $x^0\in B_{h_0}$.
The coordinate sign comparisons force every label active at $x^0$ to
have slope zero, just as they did at $c$. An active constant label
must have the least value among all constant labels, namely
$M_{h_0}=F_{h_0}(c)$. Hence
\[
 x^0\in K_{h_0}.
\]
Lemma~\ref{lem:trapping} implies that the solution through $x^0$ is
global, stays positive, and remains in $K_{h_0}$. Since $f(x)\in S$,
it also remains in $P$.

Starting at $h_0$, define
\begin{equation}\label{eq:scale-sequence}
 h_{j+1}=\min\{\sqrt{h_j},h_P\}.
\end{equation}
We claim that for every $j$ there is a finite time $T_j$ such that
$x(t)\in K_{h_j}$ for all $t\geq T_j$. The case $j=0$ holds with
$T_0=0$. Suppose it holds at $j$. Because
$h_{j+1}\leq\sqrt{h_j}$, the boxes satisfy
$B_{\sqrt{h_j}}\subset B_{h_{j+1}}$. Thus
\eqref{eq:scale-containment} yields
\[
 x(t)\in B_{h_{j+1}}\quad(t\geq T_j).
\]
The finite-entry assertion of Lemma~\ref{lem:trapping}, followed by
its invariance assertion, supplies a finite $T_{j+1}\geq T_j$ with
$x(t)\in K_{h_{j+1}}$ for all $t\geq T_{j+1}$. This proves the claim.
The argument uses containment in successive boxes; no nesting of the
sets $K_h$ is assumed.

Repeated square roots of $h_0\in(0,1)$ tend to $1$. Since $h_P<1$,
the capped sequence \eqref{eq:scale-sequence} therefore reaches $h_P$
after finitely many steps. We conclude that every solution through a
point of $P$ enters $K_{h_P}$ in finite time and stays there.

Finally set
\[
 K_P=K_{h_P}\cap(c+S).
\]
This is nonempty because it contains $c$, and it is compact and convex
as the intersection of the compact convex set $K_{h_P}$ with a closed
affine subspace. Since $K_{h_P}$ lies in the positive orthant, $K_P\subset P$.
Both $K_{h_P}$ and the affine subspace $c+S$ are forward invariant, so
$K_P$ is forward invariant. The preceding argument proves its absorbing
property. Moreover, $K_P\subset B_{h_P}$, and therefore
$\eps_P=h_P\in(0,1)$ gives
\[
 \eps_P\leq x_i(t)\leq\eps_P^{-1}
 \quad(1\leq i\leq d)
\]
for all sufficiently large $t$, with the entry time allowed to depend
on $x^0$. The scale $h_P$, and hence these bounds, was fixed before
choosing the initial point.
\end{proof}
